\section{Updating window kernels in their small parameter space}
\label{sec:window-basis}

The signed-residual filter avoids most nearby type assignments, but the
first implementation still rebuilds a matching kernel for each retained
one. Its motivating example queries 32 sectors. Each rebuild repeats
contraction, cycle construction and rational elimination on nearly the
same source. The new implementation retains the exact information computed
by the residual planner and updates the kernel in at most $d_0+2$
parameters. It also reuses source triangle potentials instead of rebuilding
the triangle graph. The independent diagram/disc certificate contract
remains the same.

\subsection{Column decompositions retain more than a residual key}

Use the fixed full-Q matrix $A$ of Section~\ref{sec:sector-windows} and
the old support $S$. During exact elimination retain coefficient vectors
alongside the column-space basis. For every nonbase column this gives
\[
                 A_j=A_Sc_j+r_j.
\]
The old planner stored only the projective direction and sign of $r_j$.
The new records also retain $c_j$ and the actual rational residual values.
The forest's original corner potential matrix $P$ is retained as sparse
integer rows. It describes signed matching triangle coordinates $Pq$;
subtracting each global vertex minimum gives the canonical nonnegative lift.
Neither record expands elementary normal pieces.

Let $K$ be a row basis of $\ker A_S$, already computed by the first sector
query. The old kernel vectors, padded with zero new coordinates, remain
kernel vectors when types are added algebraically. The following small
updates supply all remaining directions.

\begin{proposition}[Exact augmented Q kernels]
For a single zero-residual column, append the vector $(-c_j,1)$ to $K$.
For two zero-residual columns, append their two independent vectors.
For a nonzero opposite pair $a,b$, choose $\gamma>0$ with
$r_a+\gamma r_b=0$, and append
\[
                    (-c_a-\gamma c_b,1,\gamma).
\]
In each case these vectors and the padded old basis span the full augmented
matching kernel and are linearly independent.
\end{proposition}
\begin{proof}
Each displayed vector satisfies the matching equation by its column
decomposition. For zero residuals the new coordinates are independent.
For a nonzero opposite pair the quotient equation forces the two new
coordinates to be proportional with ratio $\gamma$, giving exactly one
new direction. Subtracting the appropriate displayed vector from any
augmented kernel vector leaves zero new coordinates and hence an old
kernel vector. This proves spanning, while independence follows by first
examining the new coordinates and then the independent old basis.
\end{proof}

Replacement edits must also set the replaced old coordinates to zero.
If $W$ is the augmented basis with $m\le d_0+2$ rows, impose
\[
                    \sum_{a=1}^m z_aW_{a,i}=0
\]
for each removed old index $i$. There are at most two such equations.
Their small nullspace gives the exact remaining parameter vectors; apply
them to $W$ and drop the removed coordinates. This preserves all replacement
constraints without solving a fresh wide matching system.

\subsection{Recovering the incumbent's canonical basis gauge}

The basis gauge matters for deterministic projective charts and traversal
order. The native wide elimination chooses the earliest possible pivot
columns, leaving the complementary free coordinates. Starting from the
spanning kernel rows, reduce with columns scanned from right to left.
Its pivots are those complementary free coordinates. Reverse the columns
back and order the rows by increasing free index. The result is exactly
the native nullspace basis, with an identity matrix in those free slots.

This uses thin row elimination on at most $d_0+2$ rows. The underlying
linear fact is the complementary-basis duality between a matrix's column
matroid and its kernel coordinate matroid. Equivalently, the greedy latest
independent coordinate functionals on the kernel complement the greedy
earliest independent matching columns. The normalized kernel basis is
unique once these free coordinates are fixed. Exact source comparisons
below test the resulting full basis, not merely its dimension.

\subsection{Projected triangle geometry without a new contraction}

For an updated Q basis $K'$, compute the source corner forms $PK'$ using
the retained sparse forest potentials. Equal forms at a global vertex are
identified, keeping the first physical corner as representative. This is
an exact contraction on the entire updated matching space. It can merge
more corners than the old zero-label contraction, because matching
constraints can force a nonzero-looking potential difference to vanish.
Single-form vertex groups are pure link factors and are omitted in the
canonical lift, just as before.

The canonical basis gauge provides free coordinates $F$ with
$q_{F_a}=z_a$. A corner form can therefore be expressed directly in those
physical Q coordinates, using nonzero coefficients only at the free slots.
The normalized envelope and planar producers can reuse their existing
interfaces. Potential values can be rational in this representation, but
the exact lift at an integral matching Q vector is integral by the original
source forest equations. The producer checks this integrality and validates
the full normal coordinates against the actual source.

Higher-nullity fallbacks remain complete. The projected kernel supplies
linear Q constraints
\[
                  q_j-\sum_a K'_{a,j}q_{F_a}=0\qquad(j\notin F)
\]
and triangle constraints equating each class to its vertex anchor plus
the corresponding potential difference. These describe precisely the
updated standard matching space, including one free link offset per
nontrivial vertex group. Coordinate-face extremality and the generic
support/rank tests therefore continue to apply. Nullity zero means no
non-link ray and completes without allocating another geometry kernel.

The change is implemented in \code{sector\_window\_basis.py}. It replaces
the internal per-window rebuild; the public sparse constructor, standalone
enumeration and independent coverage checkers retain their previous paths.
The source is read-only throughout the query. Positive source replay still
reconstructs geometry independently and never trusts the cached potentials,
updated basis, class identifications or planner statistics.

\subsection{Complexity and retained scope}

The residual planner pays a polynomial one-time cost for the actual column
decompositions. Per edit the matching update combines at most two new
directions with the old kernel, solves at most two equations in at most
$d_0+2$ parameters, and canonicalizes that thin basis. Triangle projection
then uses the retained sparse source rows. Wide kernel elimination and
repeated source graph contraction are removed from each candidate query.
Exact rational coefficient bits and output costs remain separately charged.

This preserves the local quasi-polynomial accounting for logarithmic
nullity and window radius from the preceding section. It neither expands
the radius-one/two window nor supplies a universal distance bound. The
default window radius remains zero unless the measured full-recognition
tradeoff justifies a later policy change. Changes to deterministic guard
counts can alter which queries complete under a finite shared work limit;
they do not turn capped local searches into absence certificates.

\subsection{Source comparisons, timings and publication checks}
All 1,404 maintained tests pass in 243.872 seconds. The focused sixteen
tests pass in 14.118 seconds. For every tested small residual window,
the updated full basis exactly equals fresh native reconstruction; the
complete ray sets, source-coordinate lifts and extremality tests agree.
An additional genuine four-tetrahedron source reaches nullity four and
checks both generic arrangement and support enumeration. A real
diagram-window test forbids more than one public full basis build and
independently replays its resulting ordinary disc proof.

The 85-source audit exactly preserves every disabled legacy result.
With radius two enabled, all verdicts equal the preceding release and
all 31 native positive witnesses remain available. The seven ordinary
window discs again receive fresh Regina confirmation. Shared-work caps
fall from 34 to 30; complete window misses rise from twenty to twenty-four.
The four additional completed cases are \code{random-8}, \code{random-14},
\code{random-41} and \code{random-135}. They still do not yield knot
verdicts. This is a guard-count improvement, not proof that every query
fits the same bit-time or wall-clock allowance.

The updated audit takes about 260.75 seconds, whereas the retained preceding
audit took about 684.44 seconds. Those are different execution records,
not a paired speedup estimate. The incumbent radius-two corpus observations
come from their frozen release data; disabled behavior is checked afresh
against the frozen implementation. All 595 captured source pins stay fixed.
The first audit attempt reached its final source and failed a research
harness tuple/list comparison; the corrected comparison and completed rerun
provide the published evidence. No interrupted run supplies final counts.

For fair native timings, the matching-rebuild, seed and recognition modules
at \code{11dde31ee} are frozen. Only their local imports are redirected
to the corresponding frozen dependency. No timed mock/patch context is
charged to either arm. Five serial interleaved four-arm rounds complete
200 measured calls and forty warmups. Seed times include full construction,
planning, actual normal queries, independent successful source replay and
serialized output. Both arms complete their configured candidate work.
The tested positive certificate digests exactly agree, as do final statuses;
implementation statistics and guard counts can differ.

\begin{samepage}
\noindent\textbf{Complete configured native window queries.}
\begin{center}
\small
\begin{tabular}{lrrrrr}
\toprule
Input & rebuild ms & update ms & old/new & old A/A & new A/A\\
\midrule
empty circle & 5.983 & 6.042 & 0.994 & 1.003 & 1.020\\
optimized positive & 60.032 & 58.637 & 0.992 & 1.011 & 0.964\\
genus one miss & 417.903 & 286.440 & 1.453 & 0.989 & 0.958\\
trefoil & 666.928 & 518.970 & 1.259 & 1.014 & 0.986\\
figure eight & 1400.078 & 967.272 & 1.453 & 1.002 & 0.994\\
\bottomrule
\end{tabular}
\end{center}

\end{samepage}

The genus-one native query drops from about 417.903 to 286.440 milliseconds,
a paired 1.453-fold gain. The trefoil and figure-eight native window queries
improve about 1.259 and 1.453 times, respectively. The motivating query
now performs one full basis build and 31 updates, retaining its 32 sector
queries, 22 emitted rays and exact disc coordinates. Early positive
controls return before the update planner and remain near parity.

\begin{samepage}
\noindent\textbf{Complete recognition with radius-two windows enabled.}
\begin{center}
\small
\begin{tabular}{lrrrrr}
\toprule
Input & rebuild ms & update ms & old/new & old A/A & new A/A\\
\midrule
empty circle & 0.020 & 0.019 & 1.036 & 0.961 & 0.948\\
optimized positive & 50.791 & 50.581 & 1.001 & 1.010 & 1.005\\
genus one miss & 427.039 & 297.222 & 1.430 & 0.981 & 1.027\\
trefoil & 696.116 & 557.498 & 1.248 & 0.990 & 1.018\\
figure eight & 1453.388 & 1020.106 & 1.431 & 1.011 & 1.001\\
\bottomrule
\end{tabular}
\end{center}

\end{samepage}

Complete enabled recognition improves about 1.430 times on the genus-one
case, 1.248 on the trefoil and 1.431 on the figure-eight. These times include
the established complete fallback where the native window misses. Every
arm reaches the expected final verdict; no partial or capped run supplies
a ratio. The improvement is relative to rebuilding the same window kernels,
not relative to disabling windows. On these small inputs, radius zero
remains substantially faster, as the preceding section records. This
optimization does not justify default activation or a universal speedup.

The empty-circle recognition control lasts only about twenty microseconds,
so its first-run A/A ratios of 0.961 and 0.948 are not interpreted as a
kernel speedup. A separate 31-round four-arm repeat completes 124 measured
calls and four warmups; its old/new ratio is 1.015, with A/A ratios 0.993
and 1.008. Both records are retained. The longer control never enters
the window planner. Other reported window-search A/A controls are retained
in the full tables.

\begin{sloppypar}
Reproduction uses \code{normal\_orbit\_research.window\_basis}, with
\code{audit --fresh-regina} or \code{benchmark --rounds 5}, from \code{fast/}.
Runtime release \code{b1741c008}, raw records, focused/full tests, the
longer control and publication checks are retained under
\code{data/window-basis-*}. Publication replay checks the frozen and current
source pins and all 31 saved positives with matching updates, projected
geometry, residual planning, coherent extraction, sector enumeration and
orbit producers disabled. Independent source-disc verification remains
the authority for positive results.
\end{sloppypar}

